\section{Forschungskonzept}
\label{sec:forschungskonzept}

Die folgenden Fragen sollen in dieser Arbeit beantwortet werden:

\begin{itemize}
    \item Entspricht Trumps Tweetverhalten dem eines klassischen Opinion-Leaders?
    \item Wer gehört zu Trumps Opinion-Followers?
    \item Wie gestaltet sich der Two-Step-Flow, wenn Trump bspw. negativ über Unternehmen tweetet?
\end{itemize}

Um den Two-Step-Flow zu visualisieren, wird der Fall Nordstrom betrachtet. Hierbei handelt es sich um eine Kaufhauskette, die sich gegen die Weiterführung der Modelinie von Trumps Tochter Ivanka entschied. Kurze Zeit später veröffentlichte Donald Trump einen Tweet darüber, dass das Unternehmen seine Tochter ungerecht behandelt habe.

Im Rahmen der Bachelorarbeit werden die folgenden methodischen Schritte umgesetzt:
\begin{enumerate}
    \item Eine Case-Study wird durchgeführt, um zu untersuchen, welche Auswirkungen der Tweet auf die Medienlandschaft hatte.
    \item Es wird analysiert, wie das Weiße Haus daraufhin kommunizierte und wie die Nachricht von den Medien aufgegriffen wurde.
    \item Der Two-Step-Flow wird visualisiert und der Agenda-Building-Prozess untersucht.
\end{enumerate}

Dabei wird die Methodik der qualitativen Inhaltsanalyse verwendet, um die Kommunikation in den sozialen Medien im Fall Nordstrom zu analysieren. Als Datensatz dienen alle relevanten Beiträge (Tweets) auf der Plattform Twitter, in denen über den Fall Nordstrom kommuniziert wird.